\section{Between the limits: closure as a function of wavenumber and frequency}
\label{sec:between}

Sections~\ref{sec:asym} and~\ref{sec:collisionless} fixed the heat flux at
the two ends: Fourier's law when collisions dominate and the
Hammett--Perkins closure when they are absent. Most applications sit in
between. This section states the problem in a form that covers every
closure met so far, and reports what is known exactly.

\subsection{One function}

For the linear problem \eqref{eq:linkin}, a single Fourier mode, and data
that lie in the retained moments, the exact relation between heat flux and
temperature is a multiplier,
\begin{equation}
  \hat q(k,\omega)=-ik\,\chi(k,\omega)\,\hat T(k,\omega) .
  \label{eq:symbol}
\end{equation}
This is the Schur complement \eqref{eq:schur} for the kinetic equation:
$\chi$ is the exact memory kernel, in Fourier--Laplace variables. Every
linear closure is a choice of approximation to $\chi$:
\begin{center}
\renewcommand{\arraystretch}{1.35}
\begin{tabularx}{\textwidth}{@{}P{4.4cm}P{4.6cm}Y@{}}
\toprule
Closure & $\chi(k,\omega)$ & Character \\
\midrule
Fourier (Chapman--Enskog) & $\kF$ & local \\
Gradient expansion, truncated & $\kF(1-Kk^2)$
  & local; negative for $k^2>1/K$ \\
Rational in $k^2$ & $\kF/(1+Kk^2)$
  & exponential kernel in space \\
Cattaneo, Grad & $\kF/(1-i\omega\tau)$
  & exponential kernel in time \\
Fokker--Planck fit (Section~\ref{sec:plasma}) & $\kF/(1+(30k\lambda_e)^{4/3})$ & fit \\
Multigroup & $\sum_jw_j/(m_j+k^2)$, $w_j\ge0$
  & sum of screened diffusions \\
Hammett--Perkins & $\sqrt{8/\pi}\;n_0\vt/|k|$
  & scale-free nonlocal \\
Harmonic interpolant & $\bigl(1/\kF+|k|/(\sqrt{8/\pi}\,n_0\vt)\bigr)^{-1}$
  & both limits by construction \\
Exact, BGK & $i\,\Phi\bigl((\omega+i\nu)/|k|\bigr)/|k|$ & see below; units $n_0=\vt=1$ \\
\bottomrule
\end{tabularx}
\end{center}
A flux limiter is not in the table because it is nonlinear: it depends on
the amplitude of the temperature gradient and has no multiplier.

Physical requirements become conditions on the function $\chi$:
causality is analyticity in the upper half of the $\omega$ plane;
dissipativity is $\operatorname{Re}\chi\ge0$ on the real axes; a real
kernel requires $\chi(-k,-\bar\omega)=\overline{\chi(k,\omega)}$; and the
two limits are $\chi\to\kF$ as $k\to0$ and $|k|\chi\to$ a constant as
$|k|\to\infty$ at $\omega=0$.

The \emph{quasi-static} closure is $\chi(k,0)$. It is the right object when
the retained dynamics is slow compared with the kinetic relaxation at that
wavenumber, and Section~\ref{sec:repo} states when that holds.

\subsection{Exact results in one velocity dimension}

Two model collision operators can be solved in closed form. Both results
are proved in the notes cited in Section~\ref{sec:repo}; the formulas were
re-evaluated here and checked against an independent Hermite--Galerkin
solve of \eqref{eq:hermitehier} with 1200 modes (agreement to
$10^{-14}$ at $\nu/|k|=1$ and $3$).

\paragraph{BGK with constant $\nu$.} Let $W(z)=\avg{1/(v-z)}$,
$R(z)=1+zW(z)$, and
\[
  \Phi(z)=\frac{z\,\bigl(1-3R(z)+z^2R(z)\bigr)}{1-R(z)+z^2R(z)} .
\]
Then every solution satisfies
$\hat q=\sgn(k)\,\Phi\bigl((\omega+i\nu)/|k|\bigr)\hat T$: the collisional
closure is the collisionless one evaluated at the complex phase velocity
$(\omega+i\nu)/|k|$. At zero frequency, with $y=\nu/|k|$,
\begin{equation}
  \hat q=-i\,\sgn(k)\,\theta(y)\,\hat T,\qquad
  \theta(y)=\frac{1+y^2}{y\,(3+y^2)}\;
  \frac{\avg{\He_3^2/(v^2+y^2)}}{\avg{\He_2^2/(v^2+y^2)}} .
  \label{eq:theta-bgk}
\end{equation}

\paragraph{Dougherty.} With $E(s)=[(1-s)e^s]^{1/y^2}$,
\begin{equation}
  \theta_1(y)=\frac2y\;\frac{\int_0^1s\,E(s)\,\dd s}{\int_0^1E(s)\,\dd s}.
  \label{eq:theta-dougherty}
\end{equation}

Both tend to $\sqrt{8/\pi}=1.5958$ as $y\to0$ (computed: $1.5949$ at
$y=0.002$ for BGK and $1.573$ at $y=0.02$ for Dougherty), and both tend to
Fourier's law as $y\to\infty$, with $\kF=3/\nu$ and $1/\nu$ respectively.
The leading deficits \eqref{eq:deficit} are $K=4/\nu^2$ and
$K=1/(3\nu^2)$ (computed at $y=50$: $3.986$ and $0.3333$).

\subsection{Comparison at matched conductivity}

Two collision operators with the same nominal $\nu$ have different
conductivities, so a fair comparison fixes $\kF$. Use the variable
$s=|k|\,\kF/(n_0\vt)$, in which Fourier's law is $\chi/\kF=1$ and the
collisionless limit is $\chi/\kF=\sqrt{8/\pi}/s$.

\begin{figure}[ht]
\centering
\begin{tikzpicture}
\begin{loglogaxis}[width=0.78\textwidth,height=7.2cm,
  xlabel={$s=|k|\,\kF/(n_0\vt)$},ylabel={$\chi/\kF$},
  xmin=0.01,xmax=100,ymin=0.01,ymax=1.5,grid=major,
  legend style={font=\footnotesize,at={(0.03,0.03)},anchor=south west}]
\addplot[black,thick] table[x=s,y=bgk]{\rundir/intermediate.dat};
\addplot[red,densely dotted,very thick] table[x=s,y=dougherty]{\rundir/intermediate.dat};
\addplot[blue,dashed,thick] table[x=s,y=harmonic]{\rundir/intermediate.dat};
\addplot[gray,thin] table[x=s,y=limits]{\rundir/intermediate.dat};
\legend{exact BGK,exact Dougherty,harmonic interpolant,the two limits}
\end{loglogaxis}
\end{tikzpicture}
\caption{Quasi-static heat-flux closure in one velocity dimension for two
collision operators at equal conductivity, with the harmonic interpolant
between Fourier's law and the Hammett--Perkins closure. Data:
\texttt{\rundir/intermediate.dat}.}
\label{fig:between}
\end{figure}

\begin{center}
\begin{tabular}{rrrr}
\toprule
$s$ & exact BGK & exact Dougherty & harmonic interpolant \\
\midrule
0.10 & 0.9956 & 0.9967 & 0.9410 \\
0.32 & 0.9595 & 0.9685 & 0.8346 \\
1.00 & 0.7570 & 0.7844 & 0.6148 \\
3.16 & 0.3884 & 0.4001 & 0.3354 \\
10.0 & 0.1466 & 0.1484 & 0.1376 \\
\bottomrule
\end{tabular}
\end{center}

Three conclusions, the first two computed here and the third quoted.
\begin{itemize}
  \item \textbf{The interpolant is biased low for these operators.} Over
    $0.01\le s\le100$ the ratio of exact to interpolated conductivity lies
    between $1.006$ and $1.23$ for BGK and between $1.006$ and $1.28$ for
    Dougherty. An interpolant that is exact at both ends is off by a
    quarter in the middle.
  \item \textbf{These two operators nearly agree at matched conductivity.}
    Their collision frequencies do not depend on speed, and the two curves
    differ by a few percent.
  \item \textbf{A speed-dependent collision frequency changes the picture.}
    For BGK with $\nu(v)=\nu_0/(1+|v|^3)$, which has the Coulomb tail, the
    note \texttt{bgk\_velocity\_dependent\_closure.tex} reports
    $\chi/\kF=0.57$ at $s\approx0.7$ and $0.13$ at $s\approx7$, far below
    the curves in Figure~\ref{fig:between}; the leading deficit coefficient
    at matched conductivity is larger by a factor of about fifteen to
    twenty; and the error of the harmonic interpolant changes sign (the
    interpolant is too high by up to $48\%$). That is the behaviour
    reported for the Landau operator \cite{Brodrick17}.
\end{itemize}
The summary is the main finding of the heat-flux manuscript: matching the
two limits and the conductivity does not determine the closure. The
coefficient $K$ in \eqref{eq:deficit} is an independent property of the
collision operator, and it is the first quantity a benchmark should test
after the conductivity itself.
